\chapter{Redes neurais que sabem aprender: acrescentamos o \nt{DOPA}}
\chaptermark{Redes que aprendem}
\label{sec:dofa}

\section{As regras do jogo}

Em todos os circuitos dos capítulos anteriores, os pesos foram ajustados pelo engenheiro. No organismo não há engenheiro. Os pesos têm de se ajustar sozinhos, durante o funcionamento, e de modo que a rede faça algo útil. É isso que se chama aprendizado.

Às regras anteriores soma-se mais uma. A sinapse muda o próprio peso sozinha e só pode saber o que acontece \emph{perto dela}: a atividade do remetente $x$, a atividade do destinatário $y$ e as substâncias que chegam até ela. Ninguém pode mandar à sinapse uma correção personalizada.

Isso exclui o principal método das redes neurais artificiais, a retropropagação do erro. Nela, o erro da saída percorre a rede no sentido inverso pelos mesmos pesos, numa fase separada depois da passagem direta, e cada peso recebe sua própria correção. Para isso são necessários fios de retorno que repitam exatamente os diretos, e um relógio comum. Nenhuma das duas coisas foi encontrada no cérebro, pelo menos na forma que têm nas redes artificiais~\cite{Lillicrap2020,WhittingtonBogacz2019}.

Vamos escrever a mudança de peso como um processo lento e contínuo:
\begin{equation}
\frac{\dd w}{\dd t} \;=\; \eta\,\Phi(\text{o que a sinapse sabe}),
\label{eq:plasticity}
\end{equation}
em que $\eta$ é a taxa de aprendizado, tão pequena que durante um disparo o peso quase não muda. O capítulo inteiro trata de qual pode ser a função $\Phi$.

\section{Onde fica o peso}
\label{sec:weightstore}

O que é uma sinapse que ``muda o próprio peso''? A conexão tem dois lados: o terminal do fio do remetente (lado \emph{pré-sináptico}) e o trecho da membrana do destinatário com receptores (lado \emph{pós-sináptico}), e entre eles uma fenda estreita. Nas conexões \nt{GLUT}, a membrana do destinatário costuma ficar numa \emph{espinha}, uma pequena saliência do ramo de entrada do destinatário. É conveniente decompor o peso da conexão em três fatores~\cite{delCastillo1954}:
\begin{empheq}[box=\keybox]{equation}
w \;=\; N_s\,p\,A_1 .
\label{eq:quantal}
\end{empheq}
Aqui $N_s$ é o número de locais de liberação entre dois neurônios, $p$ é a probabilidade de um disparo liberar uma porção de neurotransmissor num local (a mesma $p$ do capítulo~\ref{sec:code}), e $A_1$ é a resposta do destinatário a uma porção, isto é, quantos receptores a capturam. O fator $p$ mora no remetente, $A_1$ no destinatário, e $N_s$ exige os dois lados: novos locais de liberação crescem, antigos desaparecem, e isso continua a vida toda.

\emph{Quem decide é o destinatário.} Numa sinapse \nt{GLUT} típica, um dos receptores de glutamato só conduz corrente se duas condições coincidirem: o glutamato chegou do remetente e a membrana do destinatário já está excitada~\cite{Nowak1984}. É a regra de Hebb embutida numa molécula: o detector de coincidência do capítulo~\ref{sec:glutcomputer}, instalado na entrada da própria sinapse. Ao ser ativado, o receptor deixa entrar cálcio no destinatário, e o cálcio desencadeia a mudança de peso. O destinatário é o único lugar onde se veem ao mesmo tempo a entrada e a saída, por isso a decisão é tomada nele.

\emph{Qualquer lado pode executar a decisão.} A potenciação de longo prazo mais estudada ocorre no destinatário: novos receptores são inseridos na membrana~\cite{Malinow2002} e a espinha cresce~\cite{Matsuzaki2004}; cresce $A_1$. Mas o destinatário também pode mudar $p$: ele libera uma substância que atravessa a fenda de volta até o remetente (o canal de retorno do apêndice~\ref{sec:othermediators}), e este passa a liberar menos neurotransmissor~\cite{WilsonNicoll2001}. Já as mudanças de peso de curto prazo, a depressão e a facilitação do capítulo~\ref{sec:code}, são mudanças de $p$ que o próprio remetente conduz segundo sua atividade recente.

Para o engenheiro, o registrador do peso está dividido entre dois chips. A regra de aprendizado é executada pelo destinatário, mas ele pode gravar o resultado tanto em si mesmo quanto no remetente, pelo canal de retorno. Por isso, a palavra ``local'' nas regras deste capítulo não significa um lado da conexão, mas a conexão inteira.

\section{Aprender sem professor}

O mais simples que uma sinapse pode fazer é se reforçar quando remetente e destinatário estão ativos juntos:
\begin{equation}
\frac{\dd w}{\dd t} \;=\; \eta\,x\,y .
\label{eq:hebb}
\end{equation}
É a regra de Hebb~\cite{Hebb1949}. Ela é local e não exige relógio. Mas tem o mesmo problema da rede \nt{GLUT} do capítulo~\ref{sec:glutcomputer}: realimentação positiva. Um peso forte aumenta $y$, um $y$ grande reforça ainda mais o peso, e os pesos crescem sem limite. O remédio é uma realimentação negativa sobre o peso: que ele também diminua em proporção a $y^2$. Para um neurônio com entradas $\mathbf x$ e saída linear $y=\mathbf w\cdot\mathbf x$, obtemos
\begin{empheq}[box=\keybox]{equation}
\frac{\dd\mathbf w}{\dd t} \;=\; \eta\,y\,\bigl(\mathbf x - y\,\mathbf w\bigr) .
\label{eq:oja}
\end{empheq}
A regra continua local: a sinapse conhece sua entrada $x_i$, a saída $y$ e o peso $w_i$.

\begin{keyblock}
\begin{proposition}[O que a regra de Hebb encontra]
\label{prop:oja}
A regra~\eqref{eq:oja} leva o vetor de pesos a ter comprimento unitário ao longo da direção em que as entradas variam mais, ou seja, ao autovetor principal da matriz de correlação das entradas $C=\langle\mathbf x\mathbf x^{\mathsf T}\rangle$~\cite{Oja1982}. O neurônio aprende a emitir a grandeza única mais expressiva na qual suas entradas podem ser comprimidas.
\end{proposition}
\end{keyblock}

Demonstração. Tomemos a média de~\eqref{eq:oja} sobre as entradas: $\dd\mathbf w/\dd t=\eta\bigl(C\mathbf w-(\mathbf w^{\mathsf T}C\mathbf w)\,\mathbf w\bigr)$. Os pontos fixos são os autovetores de $C$ de comprimento unitário. Se $\mathbf w$ está perto do autovetor $\mathbf e_k$ com autovalor $\lambda_k$, uma pequena perturbação ao longo de outro vetor $\mathbf e_j$ cresce como $e^{\eta(\lambda_j-\lambda_k)t}$. Só é estável o ponto para o qual todos os $\lambda_j<\lambda_k$, isto é, o vetor principal.

Há também uma variante da regra de Hebb que olha para o tempo dos disparos. Se o disparo do remetente chegou $s$ milissegundos \emph{antes} do disparo do destinatário, o peso cresce $A_+e^{-s/\tau_+}$; se chegou $s$ milissegundos \emph{depois}, diminui $A_-e^{-s/\tau_-}$, com $\tau_\pm$ da ordem de vinte milissegundos. Essa regra foi medida em sinapses de neurônios \nt{GLUT}~\cite{BiPoo1998}. Ela reforça as conexões que \emph{podem ter sido a causa} da resposta e enfraquece as atrasadas (fig.~\ref{fig:stdp}). Passe muitas vezes pela rede uma mesma sequência de excitações, e nela crescerão sozinhas conexões de cada elo para o seguinte: a cadeia do capítulo~\ref{sec:glutcomputer}, montada sem engenheiro.

\begin{figure}[t]
\centering
\begin{tikzpicture}[>=Stealth,x=0.07cm,y=1.3cm]
\draw[->,gray!70!black] (-66,0) -- (68,0) node[right,black] {\small $s$};
\draw[->,gray!70!black] (0,-1.2) -- (0,1.35) node[above,black] {\small mudança de peso};
\draw[very thick,blue!60!black] plot[domain=0.5:60,samples=50] (\x,{exp(-\x/20)});
\draw[very thick,red!70!black] plot[domain=-60:-0.5,samples=50] (\x,{-0.6*exp(\x/20)});
\draw[blue!60!black,thick] (0.5,0) -- (0.5,1);
\draw[red!70!black,thick] (-0.5,0) -- (-0.5,-0.6);
\foreach \x in {-40,-20,20,40}{\draw[gray!60!black] (\x,-0.04) -- (\x,0.04);}
\node[below,font=\scriptsize] at (20,-0.04) {$20\ms$};
\node[above,font=\scriptsize] at (-20,0.04) {$-20\ms$};
\node[align=left,font=\small,blue!60!black,anchor=west] at (22,0.95) {remetente antes:\\ conexão se reforça};
\node[align=right,font=\small,red!70!black,anchor=east] at (-12,-0.95) {remetente depois:\\ conexão se enfraquece};
\end{tikzpicture}
\caption{Regra de Hebb temporal. No eixo horizontal, $s=t_{\text{dest}}-t_{\text{rem}}$, o quanto o disparo do destinatário ficou atrás do disparo do remetente; $\tau_\pm=20\ms$, $A_-=0.6A_+$. A janela, com largura de algumas dezenas de milissegundos, é da mesma ordem que a janela de coincidência~\eqref{eq:window}.}
\label{fig:stdp}
\end{figure}

Mas todas as regras desta seção ensinam o que é \emph{frequente}, e não o que é \emph{útil}. Uma sinapse que só conhece $x$ e $y$ não sabe se a ação trouxe suco ou dor. Para aprender com o resultado, a sinapse precisa de um terceiro sinal.

\section{O terceiro fator}

O terceiro sinal é trazido pelo \nt{DOPA}: os neurônios dopaminérgicos difundem um único número, o erro de predição de recompensa $\delta$~\eqref{eq:rpe} (capítulo~\ref{sec:neurontypes}). Seja a mudança de peso proporcional ao produto de três fatores: a atividade do remetente, a atividade do destinatário e $\delta$.

A dificuldade é que a recompensa não chega quando a rede escolhia a ação, mas centenas de milissegundos ou segundos depois. Quando $\delta$ chega, o produto $xy$ já é zero há muito tempo, e a sinapse precisa de uma memória de que participou. A memória mais simples é o nosso conhecido circuito RC:
\begin{empheq}[box=\keybox]{equation}
\tau_e\,\frac{\dd e}{\dd t} \;=\; -e \;+\; x\,y ,
\qquad
\frac{\dd w}{\dd t} \;=\; \eta\,\delta(t)\,e(t) .
\label{eq:threefactor}
\end{empheq}
A grandeza $e$ se chama \emph{traço de elegibilidade}~\cite{Fremaux2016}. A atividade conjunta deixa na sinapse uma marca que se desfaz no tempo $\tau_e$. Se nesse tempo chega dopamina, o peso muda com o sinal de $\delta$; se não chega, a marca some sem deixar rastro (fig.~\ref{fig:trace}). Em sinapses \nt{GLUT} foi medido: a dopamina que chega um ou dois segundos depois da atividade conjunta a converte em reforço da conexão, e a que chega mais tarde, não~\cite{Yagishita2014}. Janelas de plasticidade de segundos também foram encontradas onde o terceiro sinal não é a dopamina, mas uma forte excitação do próprio destinatário~\cite{Bittner2017}.

\begin{figure}[t]
\centering
\begin{tikzpicture}[>=Stealth,x=7cm,y=1cm]
\fill[orange!12] (0.7,-0.2) rectangle (0.8,5.2);
% rows: x*y, delta, traces, weights
\foreach \y/\lab in {4.4/{$x\,y$},3.1/{$\delta$},1.4/{traço $e$},0/{peso $w$}}{
  \draw[gray!70!black] (0,\y) -- (1.42,\y);
  \node[left,font=\small] at (-0.02,\y+0.3) {\lab};}
\draw[very thick,blue!60!black] (0,4.4) -- (0,4.9) -- (0.2,4.9) -- (0.2,4.4);
\node[right,font=\scriptsize,blue!60!black] at (0.21,4.8) {remetente e destinatário ativos juntos};
\draw[very thick,orange!85!black] (0,3.1) -- (0.7,3.1) -- (0.7,3.7) -- (0.8,3.7) -- (0.8,3.1) -- (1.4,3.1);
\node[right,font=\scriptsize,orange!85!black] at (0.81,3.6) {chegou dopamina};
% traces, each in units of its own maximum
\draw[very thick,blue!60!black] plot[domain=0:0.2,samples=20] (\x,{1.4+1.1*(1-exp(-\x))/(1-exp(-0.2))})
  plot[domain=0.2:1.4,samples=60] (\x,{1.4+1.1*exp(-(\x-0.2))});
\draw[thick,densely dashed,gray!50!black] plot[domain=0:0.2,samples=30] (\x,{1.4+1.1*(1-exp(-\x/0.05))/(1-exp(-4))})
  plot[domain=0.2:1.4,samples=80] (\x,{1.4+1.1*exp(-(\x-0.2)/0.05)});
\node[right,font=\scriptsize,blue!60!black] at (1.0,2.15) {$\tau_e=1\,\text{s}$};
\node[right,font=\scriptsize,gray!40!black] at (0.3,1.65) {$\tau_e=50\ms$};
% weights
\draw[very thick,blue!60!black] (0,0.25) -- (0.7,0.25) -- (0.8,0.8) -- (1.4,0.8) node[right,font=\scriptsize] {$\tau_e=1\,\text{s}$: peso cresceu};
\draw[thick,densely dashed,gray!50!black] (0,0.2) -- (1.4,0.2) node[right,font=\scriptsize,gray!40!black] {$\tau_e=50\ms$: sem mudança};
% time axis marks
\foreach \x/\l in {0/0,0.2/0.2,0.7/0.7,1.4/1.4}{\draw[gray!60!black] (\x,-0.05) -- (\x,0.05) node[below=3pt,font=\scriptsize] {\l\,s};}
\draw[<->,thick,gray!50!black] (0.2,5.35) -- node[above,font=\scriptsize,black] {atraso da recompensa $0.5\,\text{s}$} (0.7,5.35);
\end{tikzpicture}
\caption{Traço de elegibilidade pela regra~\eqref{eq:threefactor}. A atividade conjunta dura $0.2\,\text{s}$, e a dopamina chega meio segundo após seu fim. Os traços são mostrados em frações do próprio máximo. O traço lento ($\tau_e=1\,\text{s}$) ainda está vivo nesse momento; o rápido ($\tau_e=50\ms$) já se desfez há muito.}
\label{fig:trace}
\end{figure}

\begin{keyblock}
\begin{proposition}[Aprendizado por sinal de difusão]
\label{prop:threefactor}
A regra~\eqref{eq:threefactor} muda só as sinapses que participaram do funcionamento da rede nos últimos $\tau_e$, no sentido dado pelo sinal comum $\delta$. O sinal de difusão, junto com o traço local, produz uma correção endereçada. Um atraso entre ação e resultado é admissível enquanto não passar de $\tau_e$.
\end{proposition}
\end{keyblock}

\section{Por que funciona: o ruído como busca}
\label{sec:dofagradient}

Por que a regra~\eqref{eq:threefactor} leva a uma melhora? A resposta vem do ruído do capítulo~\ref{sec:code}. Seja a saída do neurônio $y=f(u)+\xi$, com $u=\sum_jw_jx_j$, e $\xi$ uma flutuação aleatória de média zero e variância $\sigma^2$. A recompensa $R$ depende de $y$. Tomemos como traço não simplesmente $x_jy$, mas sua parte aleatória $x_j\xi$, e como terceiro fator o erro $\delta=R-\bar R$, em que $\bar R$ é a recompensa esperada. Com ruído pequeno, $R\approx R\bigl(f(u)\bigr)+R'\,\xi$, portanto
\begin{empheq}[box=\keybox]{equation}
\bigl\langle \delta\,\xi\,x_j \bigr\rangle \;=\; \sigma^2\,R'\,x_j \;=\; \frac{\sigma^2}{f'(u)}\,\frac{\partial\langle R\rangle}{\partial w_j} .
\label{eq:perturbation}
\end{empheq}
O passo médio segue o gradiente da recompensa esperada~\cite{Williams1992}. Ninguém calculou derivadas: a rede se desviou ao acaso, a dopamina informou se melhorou, e as sinapses que participaram se deslocaram no sentido vantajoso (fig.~\ref{fig:perturbation}). O ruído, que no capítulo~\ref{sec:code} atrapalhava a transmissão de mensagens, aqui se torna busca.

\begin{figure}[t]
\centering
\begin{tikzpicture}[>=Stealth,x=2cm,y=3cm]
\draw[->,gray!70!black] (0,0) -- (4.2,0) node[below left,font=\small,black] {saída $y$};
\draw[->,gray!70!black] (0,0) -- (0,1.2) node[left,font=\small,black] {recompensa $R$};
% reward hill
\draw[very thick,blue!60!black] plot[domain=0:4,samples=60] (\x,{max(0,1-((\x-3)/2.2)^2)});
\node[font=\scriptsize,blue!60!black,anchor=south] at (3,1.02) {melhor saída};
% noise around f(u)
\draw[thick,gray!60!black] plot[domain=0.6:2.4,samples=50] (\x,{0.28*exp(-((\x-1.5)/0.3)^2/2)});
\draw[densely dashed,gray!60!black] (1.5,0) -- (1.5,0.535);
\node[below,font=\small] at (1.5,0) {$f(u)$};
\node[font=\scriptsize,gray!50!black] at (1.5,-0.2) {flutuação $\xi$};
% expected reward
\draw[densely dotted,gray!60!black] (0.5,0.535) node[left,font=\scriptsize,black] {$\bar R$} -- (2.1,0.535);
% two random deviations
\fill[green!45!black] (2.1,0.833) circle (0.035);
\draw[very thick,green!45!black] (2.1,0.535) -- (2.1,0.833);
\node[font=\scriptsize,green!45!black,anchor=west,align=left] at (2.18,0.6) {$\xi>0$: melhor,\\ $\delta>0$};
\fill[red!70!black] (0.9,0.089) circle (0.035);
\draw[very thick,red!70!black] (0.9,0.535) -- (0.9,0.089);
\node[font=\scriptsize,red!70!black,anchor=east,align=right] at (0.83,0.27) {$\xi<0$: pior,\\ $\delta<0$};
% resulting step
\draw[->,very thick,orange!85!black] (1.55,0.4) -- (1.95,0.4) node[right,font=\scriptsize,black] {passo};
\end{tikzpicture}
\caption{O ruído como busca~\eqref{eq:perturbation}. A saída do neurônio flutua em torno de $f(u)$. Um desvio aleatório para a direita (verde) deu mais que o esperado, $\delta>0$; para a esquerda (vermelho), menos, $\delta<0$. Nos dois casos o produto $\delta\,\xi$ é positivo, e o peso se desloca para a direita, onde a recompensa cresce. Em média, o passo é proporcional à inclinação $R'$: no topo da colina, desvios para os dois lados são igualmente ruins, e os passos se anulam.}
\label{fig:perturbation}
\end{figure}

\begin{keyblock}
\begin{proposition}[Descida de gradiente sem fios de retorno]
\label{prop:gradient}
Se a rede tem desvios aleatórios de atividade e o sinal de difusão é igual ao erro de predição de recompensa, com média zero em cada situação, a regra~\eqref{eq:threefactor} muda, em média, os pesos segundo o gradiente da recompensa esperada. Para isso não são necessários fios de retorno nem uma fase separada de aprendizado.
\end{proposition}
\end{keyblock}

Agora fica claro por que a dopamina não informa a recompensa, mas o \emph{erro} de sua predição. Se o terceiro fator fosse a própria recompensa $R\ge0$, a média em~\eqref{eq:perturbation} não mudaria, mas surgiriam dois problemas. Primeiro, à parte útil se somaria o termo $\bar R\,\xi x_j$: zero em média, mas um ruído forte. Segundo, e pior, uma sinapse real não sabe separar no traço $x_jy$ a parte aleatória da não aleatória. Com as entradas dadas, $x_jf(u)$ é o mesmo número de uma tentativa para outra; se $\delta$ tem média zero nessa situação, essa parte do traço não muda nada, mas com $\delta\ge0$ ela reforça, vez após vez, tudo o que a rede fez, o certo e o errado. Por isso $\bar R$ deve ser a expectativa \emph{na situação dada}, e não a média de toda a vida. Calculá-la é tarefa do crítico, de que falamos abaixo.

Verificamos isso num modelo. Dois grupos de neurônios \nt{GLUT} escolhem uma de duas ações por inibição mútua, como na fig.~\ref{fig:wta}; os pesos do sinal ``começar'' para os grupos são inicialmente iguais, e o ruído decide quem vence. A ação $A$ traz suco com probabilidade $0.8$, a ação $B$ com probabilidade $0.2$, e o suco chega meio segundo depois da escolha. Com $\tau_e=1\,\text{s}$ e $\delta=R-\bar R$, a fração de escolhas de $A$ em quinhentas tentativas subiu de $0.65$--$0.8$ na primeira centena para $0.96$--$1.0$ na última, e os pesos ficaram entre $0.85$ e $1.35$. Com $\tau_e=50\ms$, menor que o atraso da recompensa, a rede não aprendeu nada: a fração de escolhas de $A$ ficou em torno da metade. Com $\delta=R$, sem subtrair a expectativa, a rede também passou a escolher $A$, mas o peso do vencedor cresceu mil vezes nas mesmas quinhentas tentativas e continuou crescendo; e com o mesmo $\delta=R$ e uma taxa de aprendizado dez vezes maior, a rede, numa de três execuções, fixou para sempre sua primeira escolha aleatória, a pior.

\section{O preço de um único fio}

O sinal $\delta$ é um só para todos, e o ruído que ele avalia é a soma das flutuações de todos os $N$ neurônios que influem no resultado. Cada sinapse vê em $\delta$ sua parte útil e o ruído alheio de $N-1$ vizinhos. Para separar o que é seu, é preciso fazer a média sobre muitas tentativas, e o tempo de aprendizado cresce proporcionalmente a $N$~\cite{Werfel2005}. A retropropagação não cobra esse preço.

\begin{keyblock}
\begin{proposition}[O preço da difusão]
\label{prop:broadcastcost}
O tempo de aprendizado por um único sinal comum cresce proporcionalmente ao número de neurônios cujo ruído influi no resultado. Esse aprendizado serve para circuitos pequenos, que escolhem entre poucas opções, mas não para ajustar bilhões de pesos.
\end{proposition}
\end{keyblock}

O cérebro, ao que parece, divide o trabalho assim: as saídas dos neurônios \nt{DOPA} vão sobretudo às regiões que escolhem ações (capítulo~\ref{sec:neurontypes}), e a enorme rede \nt{GLUT} do córtex, onde está a esmagadora maioria dos pesos, aprende principalmente por regras locais, sem professor externo. Antes, elas eram reduzidas às regras de Hebb~\eqref{eq:oja}. Hoje há cada vez mais dados de que o córtex tem um objetivo próprio, \emph{prever} sua próxima entrada, e então o erro é calculado no próprio local, como a diferença entre o previsto e o que chegou~\cite{Keller2018}: o professor está embutido na própria rede. Janelas de plasticidade de segundos, em que o terceiro sinal é uma forte excitação do próprio destinatário~\cite{Bittner2017}, mostram que um sinal de erro local nas sinapses de fato existe. Até que ponto essa é uma regra geral do córtex é uma hipótese.

\section{De onde vem o erro: o crítico em tempo contínuo}

Resta a pergunta de como os próprios neurônios \nt{DOPA} calculam $\delta$. Nos programas de aprendizado por reforço, o erro de predição é calculado em passos $t,t+1,\dots$ No organismo não há passos, por isso vamos escrevê-lo em tempo contínuo~\cite{Doya2000}. Seja $r(t)$ o fluxo de recompensa e $V(t)$ a expectativa de toda a recompensa futura, em que a mais distante vale menos:
\begin{equation}
V(t) \;=\; \int_t^{\infty} e^{-(s-t)/\tau_\gamma}\,r(s)\,\dd s .
\end{equation}
Derivando, obtemos $\dd V/\dd t=V/\tau_\gamma-r$. O resíduo dessa igualdade é justamente o erro de predição:
\begin{empheq}[box=\keybox]{equation}
\delta(t) \;=\; r(t) \;+\; \frac{\dd V}{\dd t} \;-\; \frac{V}{\tau_\gamma} .
\label{eq:ctd}
\end{empheq}
A constante $\tau_\gamma$ é o horizonte de planejamento, análogo contínuo do coeficiente $\gamma$; pela hipótese, é definida pelo \nt{SERO} (seção~\ref{sec:horizon}), e então~\eqref{eq:patience} é a regra de quanto vale a pena esperar. Vamos verificar três casos (fig.~\ref{fig:ctdcases}). Acendeu-se uma luz que promete suco: $V$ subiu de repente, $\dd V/\dd t$ é grande, pico. Chegou o suco prometido: $r>0$, mas a expectativa se esgotou no mesmo instante, $\dd V/\dd t\approx-r$, e não há pico. O suco não veio: a expectativa sumiu sem recompensa, $\dd V/\dd t<0$, queda.

\begin{figure}[t]
\centering
\begin{tikzpicture}[>=Stealth,x=1.15cm,y=1cm]
\foreach \col/\ttl/\lampon/\juice in {0/{suco inesperado}/0/1, 1/{suco prometido}/1/1, 2/{o suco não veio}/1/0}{
 \begin{scope}[xshift=\col*4.6cm]
  \node[font=\small] at (1.5,3.55) {\ttl};
  \foreach \yy in {2.4,1.2,0}{\draw[gray!60!black] (0,\yy) -- (3.1,\yy);}
  % reward r
  \ifnum\juice=1
    \fill[green!45!black] (1.95,2.4) rectangle (2.05,2.85);
    \pic[scale=0.55] at (2.0,3.1) {juice};
  \else
    \pic[scale=0.55] at (2.0,3.1) {nojuice};
  \fi
  % expectation V and error delta
  \ifnum\lampon=1
    \pic[scale=0.55] at (1.0,3.1) {lamp};
    \draw[very thick,blue!60!black] (0,1.2) -- (1,1.2) -- (1,1.45)
      plot[domain=1:2,samples=20] (\x,{1.2+0.25*exp((\x-1)/1.5)}) -- (2,{1.2+0.25*exp(1/1.5)}) -- (2,1.2) -- (3.1,1.2);
    \draw[very thick,red!70!black] (0,0) -- (0.95,0) -- (0.95,0.5) -- (1.05,0.5) -- (1.05,0) -- (3.1,0);
    \ifnum\juice=0
      \draw[very thick,red!70!black] (1.95,0) -- (1.95,-0.45) -- (2.05,-0.45) -- (2.05,0);
      \node[font=\scriptsize,red!70!black,anchor=west] at (2.1,-0.35) {queda};
    \else
      \node[font=\scriptsize,gray!50!black,anchor=north,align=center] at (2.0,-0.08) {$r$ e a queda de $V$\\ se anulam};
    \fi
  \else
    \draw[very thick,blue!60!black] (0,1.2) -- (3.1,1.2);
    \draw[very thick,red!70!black] (0,0) -- (1.95,0) -- (1.95,0.5) -- (2.05,0.5) -- (2.05,0) -- (3.1,0);
  \fi
  \ifnum\col=0
    \node[font=\small,anchor=east] at (-0.1,2.55) {$r$};
    \node[font=\small,anchor=east] at (-0.1,1.35) {$V$};
    \node[font=\small,anchor=east] at (-0.1,0.1) {$\delta$};
  \fi
 \end{scope}}
\end{tikzpicture}
\caption{Três casos para o erro~\eqref{eq:ctd}; de cima para baixo, recompensa $r$, expectativa $V$ e erro $\delta$. À esquerda não há expectativa, e o suco é inteiramente uma surpresa. No meio, a luz eleva $V$ de repente: é um salto de $\dd V/\dd t$, e o pico de $\delta$ cai na luz. Depois, $V$ cresce exatamente como o horizonte exige, e $\delta=0$; no instante do suco, a recompensa $r$ e a queda de $V$ se anulam. À direita, $V$ cai sem recompensa: queda. São o mesmo pico, a mesma transferência e a mesma queda da fig.~\ref{fig:rpe}, mas agora se vê de onde vêm.}
\label{fig:ctdcases}
\end{figure}

O que é preciso, dos circuitos que já temos, para~\eqref{eq:ctd}? Três coisas.

\emph{A própria estimativa $V$.} Ela é emitida por um grupo de neurônios \nt{GLUT}, o crítico, cujos pesos aprendem pela mesma regra~\eqref{eq:threefactor} com o mesmo $\delta$: se $\delta>0$, a expectativa estava subestimada, e as conexões ativas logo antes se reforçam.

\emph{A derivada $\dd V/\dd t$.} Ela é calculada por qualquer circuito que responda às variações, e não ao nível: excitação direta junto com inibição atrasada por um intermediário \nt{GABA}, como no detector de direção do capítulo~\ref{sec:gaba}, ou uma sinapse com depressão~\eqref{eq:depression} do capítulo~\ref{sec:code}.

\emph{O sinal num único fio.} O erro pode ser positivo ou negativo, mas a taxa de disparos só pode ser positiva. A solução é a mesma do neurônio \nt{SERO} do capítulo~\ref{sec:serocomputer}: o neurônio \nt{DOPA} é um marca-passo. Seus poucos disparos regulares por segundo significam $\delta=0$; uma rajada, $\delta>0$; uma pausa, $\delta<0$. Para o erro negativo sobra menos espaço: a taxa não pode cair abaixo de zero. Nos neurônios \nt{DOPA} reais, as quedas são de fato menores que os picos~\cite{Bayer2005}.

Que a dopamina se comporta como $\delta$ foi medido. Como exatamente o cérebro calcula $\dd V/\dd t$ é uma hipótese.

\section{Ator e crítico}

Juntemos tudo (fig.~\ref{fig:actorcritic}). Os neurônios de contexto excitam grupos de ação, que escolhem um vencedor via \nt{GABA} (proposição~\ref{prop:wta}), o \emph{ator}, e o crítico, que emite $V$~\cite{SuttonBarto2018}. O neurônio \nt{DOPA} recebe a recompensa $r$ e $V$, calcula~\eqref{eq:ctd} e difunde $\delta$ a todas as sinapses do ator e do crítico.

\begin{figure}[t]
\centering
\begin{tikzpicture}[>=Stealth,
  nrn/.style={circle,draw,very thick,fill=blue!6,minimum size=0.9cm,inner sep=0pt,font=\small},
  gab/.style={circle,draw,very thick,fill=red!10,minimum size=0.6cm,inner sep=0pt},
  dofa/.style={circle,draw,very thick,double,double distance=1.2pt,fill=orange!15,minimum size=1.35cm,inner sep=0pt,font=\scriptsize},
  inh/.style={-{Bar[width=7pt]},thick,red!70!black},
  bc/.style={->,thick,densely dashed,orange!85!black}]
\node[nrn] (S1) at (0,2.4) {$x_1$};
\node[nrn] (S2) at (0,1.2) {$x_2$};
\node[nrn] (S3) at (0,0) {$x_3$};
\node[left=0.15cm,align=right,font=\small] at (-0.5,1.2) {contexto};
\node[nrn] (A) at (4,2.6) {$A$};
\node[nrn] (B) at (4,1.0) {$B$};
\node[gab] (G) at (5.2,1.8) {};
\draw[->,thick] (A) -- (G); \draw[->,thick] (B) -- (G);
\draw[inh] (G) to[bend left=35] (A);
\draw[inh] (G) to[bend right=35] (B);
\node[above,font=\small] at (4.6,3.05) {ator: escolha da ação};
\node[nrn] (V) at (4,-1.1) {$V$};
\node[below,font=\small] at (4,-1.6) {crítico};
\foreach \s in {S1,S2,S3}{\foreach \t in {A,B,V}{\draw[->,gray!60!black] (\s) -- (\t);}}
\node[dofa] (D) at (8.0,-1.1) {\nt{DOPA}};
\draw[->,thick] (V) -- node[above,font=\scriptsize] {$V,\ \dd V/\dd t$} (D);
\draw[->,thick,green!45!black] (10.0,-1.1) node[right,black,font=\small] {recompensa $r$} -- (D);
\pic[scale=1.1] at (10.6,-0.5) {juice};
\draw[bc] (D.north) -- (8.0,4.0) -- node[above,font=\small,black] {$\delta$ a todas as sinapses do ator e do crítico} (2.2,4.0) -- (2.2,3.0);
\draw[bc] (D.south) -- (8.0,-2.4) -- (2.2,-2.4) -- (2.2,-0.6);
\end{tikzpicture}
\caption{Uma rede que aprende a escolher ações. Setas cinza: sinapses \nt{GLUT} com pesos variáveis; círculo vermelho: neurônio \nt{GABA}; círculo duplo: marca-passo \nt{DOPA}, que calcula $\delta$ por~\eqref{eq:ctd}; setas tracejadas: difusão de $\delta$.}
\label{fig:actorcritic}
\end{figure}

O sinal da ação do neurotransmissor é escolhido pelo destinatário (proposição~\ref{prop:receptor}), e na região de escolha de ações parte dos neurônios tem receptores de dopamina de uma família, e parte, de outra. Em experimentos com fatias, a dopamina junto com a atividade conjunta reforça as sinapses nos primeiros e as enfraquece nos segundos~\cite{Shen2008}; no modelo, isso significa que nos segundos o sinal de $\delta$ em~\eqref{eq:threefactor} está invertido. Os primeiros aprendem o que \emph{vale a pena fazer}; os segundos, o que \emph{não vale a pena fazer}.

\section{Resumo}

\begin{table}[t]
\centering
\footnotesize
\setlength{\tabcolsep}{4pt}
\renewcommand{\arraystretch}{1.15}
\begin{tabular}{>{\raggedright}p{2.6cm}>{\raggedright}p{2.7cm}>{\raggedright}p{3.1cm}>{\raggedright\arraybackslash}p{5.0cm}}
\toprule
Regra & O que a sinapse sabe & O que a rede aprende & O que se sabe \\
\midrule
Hebb por taxas~\eqref{eq:oja} & $x$, $y$, o próprio peso & comprimir as entradas: direções principais & reforço com atividade conjunta medido; o mecanismo de limitação dos pesos é objeto de pesquisa \\
Hebb temporal & ordem dos disparos de $x$ e $y$ dentro de $\sim20\ms$ & conexões causais, sequências & medido em sinapses \nt{GLUT} \\
três fatores~\eqref{eq:threefactor} & $x$, $y$, traço $e$, sinal comum $\delta$ & ações que trazem recompensa & efeito da dopamina nos segundos após a atividade medido; como principal mecanismo do aprendizado de escolha, hipótese bem fundamentada \\
crítico~\eqref{eq:ctd} & o mesmo & expectativa de recompensa $V$ & semelhança da dopamina com $\delta$ medida; o circuito que calcula $\dd V/\dd t$ é hipótese \\
retropropagação do erro & correção personalizada para cada peso & tudo, mas exige fios de retorno e relógio & não encontrada no cérebro nessa forma; procuram-se aproximações \\
\bottomrule
\end{tabular}
\caption{Regras de aprendizado numa rede de neurônios \nt{GLUT}, \nt{GABA} e \nt{DOPA}.}
\label{tab:learning}
\end{table}

As regras de aprendizado estão reunidas na tab.~\ref{tab:learning}. O \nt{GLUT} transporta os dados, o \nt{GABA} controla o fluxo, e o \nt{DOPA} decide quais mudanças na rede ficam.

\section{Exercícios}

\begin{exercise}[\lvlA]
\label{ex:06:1}
\emph{Quanto dura o traço.} A atividade conjunta $xy=1$ dura $T_a=0.2\,\text{s}$, a partir de $e=0$, e a dopamina chega $d=0.5\,\text{s}$ após seu fim. (a)~Quantas vezes o traço~\eqref{eq:threefactor} nesse instante é maior com $\tau_e=1\,\text{s}$ do que com $\tau_e=50\ms$? (b)~Com que $\tau_e$ ele é máximo?
\end{exercise}

\begin{exercise}[\lvlA]
\label{ex:06:2}
\emph{Deriva da regra de Hebb.} Remetente e destinatário emitem trens de Poisson independentes de $10$ disparos por segundo, e cada par de disparos muda o peso pela janela da fig.~\ref{fig:stdp} com $\tau_\pm=20\ms$, $A_-=0.6A_+$. Quantos $A_+$ o peso cresce em uma hora de atividade totalmente aleatória? Com que $\tau_-$ a deriva desaparece?
\end{exercise}

\begin{exercise}[\lvlB]
\label{ex:06:3}
\emph{Correlações, não covariâncias.} A regra da proposição~\ref{prop:oja} encontra o autovetor principal de $C=\langle\mathbf x\mathbf x^{\mathsf T}\rangle$. As taxas são positivas: $\mathbf x=\mathbf m+\mathbf n$, $\mathbf m=(\mu,0)$, covariância do ruído $\operatorname{diag}(s_1^2,s_2^2)$. Sob que condição o neurônio aprende $\mathbf e_1$? O que aprende com $\mu=1$, $s_1=0.1$, $s_2=0.5$?
\end{exercise}

\begin{exercise}[\lvlA]
\label{ex:06:4}
\emph{O horizonte no pico.} A luz acende num instante imprevisível, e um suco de tamanho $R$ chega após $L=1\,\text{s}$. Mostre que, para a expectativa aprendida por~\eqref{eq:ctd}, entre a luz e o suco $\delta\equiv0$, e encontre a área do pico na luz com $\tau_\gamma=2\,\text{s}$ e $\tau_\gamma=4\,\text{s}$.
\end{exercise}

\begin{exercise}[\lvlB]
\label{ex:06:5}
\emph{De onde vem o preço da difusão.} $N$ neurônios flutuam independentemente, $\xi_i\sim\mathcal N(0,\sigma^2)$, o erro é $\delta=a\sum_i\xi_i$, e a sinapse $j$ estima o gradiente como $\delta\,\xi_jx_j$. Encontre a relação sinal-ruído de uma tentativa. Um neurônio aprende em $100$ tentativas de $5\,\text{s}$: quanto leva o aprendizado com $N=10^4$ e $N=10^{10}$?
\end{exercise}

\begin{exercise}[\lvlB]
\label{ex:06:6}
\emph{Projeto: circuito que calcula $\delta$.} O neurônio \nt{DOPA} recebe $r$, $V$ com peso $k_1$ e uma cópia de $V$ suavizada com $\tau_d$, com peso $-k_2$. (a)~Escolha $k_1$, $k_2$ para que, com $V$ lento, a saída seja~\eqref{eq:ctd}. (b)~Com $V=V_0e^{t/\tau_\gamma}$, o erro verdadeiro é $\delta=0$. Com que $\tau_d$ o erro espúrio do circuito não passa de $5\,\%$ de $V/\tau_\gamma$, se $\tau_\gamma=2\,\text{s}$?
\end{exercise}

\begin{exercise}[\lvlB]
\label{ex:06:7}
\emph{Simulação: atraso máximo da recompensa.} Acrescente um atraso de recompensa $d$ em \texttt{run()} numa cópia de \texttt{models/\allowbreak dofa\_learning.py}. Ache a fração de escolhas de $A$ nas últimas $100$ de $500$ tentativas com $d=0.5\,\text{s}$ e $\tau_e=0.05$, $0.2$, $1$, $4\,\text{s}$, e com $\tau_e=1\,\text{s}$, $d=3\,\text{s}$. Com que fração do traço vivo até a recompensa (exercício~\ref{ex:06:1}) o aprendizado some?
\end{exercise}

\begin{exercise}[\lvlC]
\label{ex:06:8}
\emph{Dá para escapar do preço da difusão?} $N$ neurônios: $y_i=w_i+\xi_i$, $\xi_i\sim\mathcal N(0,\sigma^2)$, recompensa $R=-\sum_i(y_i-w_i^*)^2$, regra $\Delta w_i=\eta\,(R-\langle R\rangle)\,\xi_i$. Em quantas tentativas, com o melhor $\eta$, o erro $\sum_i(w_i-w_i^*)^2$ cai $e$ vezes? E se só $m$ neurônios aleatórios de $N$ flutuarem? A busca esparsa ajuda?
\end{exercise}
